% !!! GENERATED FILE - DO NOT EDIT !!!
% Produced from paper/sections/06_positioning.html by paper/assemble_arxiv.py
% (rules in arxiv/_CONVERSION_SPEC.md). Edit the .html source + meta.py, then
% rerun `python paper/build.py` (or `python paper/assemble_arxiv.py`).

\section{Is the signal real? Naive baselines}\label{sec:baselines}

A detector that fires on the right runs is only impressive if a dumb
detector cannot do the same. So we measured the dumb ones, scored by the
same lift and false-alarm from \cref{sec:bench} (lift = precision minus the 76.2\%
base fail rate; false-alarm = how often the detector fires on a run that
actually passed). A good detector has high lift and low false-alarm.

\begin{table}[t]
  \centering
  \caption{\cref{tab:baselines}. Naive baselines vs. the shipped detector
(Toolathlon, third-party scored). Read the \emph{lift} column: it is precision
minus the 76.2\% base, so it strips out the credit a detector gets for free on a
benchmark where failure is the default.}
  \label{tab:baselines}
  \begin{tabular}{lrrrr}
    \toprule
    Detector & Fire & Precision & Lift & False-alarm \\
    \midrule
    baseline-always-fail & 100\% & 76.2\% & +0.0pp & 100\% \\
    naive-loose & 32.3\% & 81.1\% & +4.9pp & 25.6\% \\
    naive-any-error & 19.5\% & 81.4\% & +5.2pp & 15.2\% \\
    tight-no-recovery & 2.3\% & 93.0\% & +16.9pp & 0.7\% \\
    terminal\_error (shipped) & 1.2\% & 95.0\% & +18.8pp & 0.24\% \\
    \bottomrule
  \end{tabular}
\end{table}

(``Fire'' is how often the detector triggers at all.)

\textbf{The recall mirage.} Recall is the wrong scoreboard here. The baseline that prints
``fail'' on every run posts 100\% recall --- it catches every failure because it flags everything
--- yet its lift is +0.0pp and it false-alarms 100\% of the time. Score by lift and false-alarm
instead and it falls to the bottom where it belongs. Our detectors make the opposite trade (high
precision, low false-alarm, low recall) because a verdict is meant to be \emph{acted on}, and
\cref{sec:giveup} is what that looks like: you only give up on a run you are sure is doomed.

\textbf{The tight grammar is load-bearing.} You might think reading the structured error
envelope more loosely would catch more failures. It does not: loosening buys $\sim$5pp of precision but
pushes false-alarm from 0.24\% to 25.6\%--69.4\% ($\approx$105$\times$--285$\times$ worse). You cannot buy
recall by reading the bytes you already have more loosely; the only way to catch more is an
\emph{independent} witness --- exactly the Tier-B re-read of \cref{sec:payoff}. The per-model
grid (\cref{fig:permodel}) shows the precision holding across every
model that fires.

\begin{figure}[t]
  \centering
  \includegraphics[width=\linewidth]{fig2_per_model_grid}
  \caption{\cref{fig:permodel}. Model by model (weakest to strongest): precision stays high all the way up the ladder; only the fire-rate falls away on the strongest models. The only honest way to catch more failures is a new source of evidence, not a looser reading of the bytes already in hand.}
  \label{fig:permodel}
\end{figure}

\section{Positioning}\label{sec:positioning}

Three kinds of system sit near this work. The first does the same job by a cheaper-looking trick --- re-run the agent until a check passes --- and works only for one agent. The second does the same job --- catch a silent failure --- by a different method (a trained classifier). The third does a different job entirely. We take each on its own terms.

\subsection{The re-run wrapper (and why a fleet defeats it)}\label{sec:positioning-wrapper}

The most common reaction is that this looks unnecessary: \emph{if you do not trust an agent's
``done,'' just check the result and re-run when the check fails.} The loop is correct for
a single agent and we make no claim against it --- only against the dismissal it smuggles in, that
single-agent reliability is therefore trivial. Two parts of the loop are real work the phrase hides:
the \emph{check} must be one the agent cannot forge (a git-ancestry, registry-first ship oracle under
a closed vocabulary, not ``re-read the result,'' which falls into the same mirror-verifier
trap), and the \emph{loop} must sense fast enough to gate every turn (here the in-flight path needed
a rewrite from a $\sim$250~ms Python hook to a $\sim$10~ms native one). Neither is a few lines of
glue.

But the \emph{headline} claim --- the place a re-run wrapper provably cannot follow --- is
the concurrent one, and we concede the stronger objection first. If the shared store offered a
transactional primitive (a serializable transaction, a compare-and-swap, a conditional write), a
wrapper around even a lone agent could fuse check and write into one atomic step and close the race
with no referee. Our claim is scoped to the regime that holds for agents in the wild: they act through
high-level tool APIs (cancel a reservation, file a ticket) that expose only unconditional writes, so no
such primitive is available to wrap. There the TOCTOU argument bites --- A's read and write
are separate non-atomic calls, so no number of re-runs makes two calls one. The number settles it:
\cref{sec:payoff-race}'s arbiter prevents all 6 live would-be clobbers across its 8 racing
pairs and a re-run wrapper prevents none, because each agent's own check passed. (The arbiter is a lock service supplying the
serialization the store does not, lifted \emph{outside} any single agent --- which is why it
works where a wrapper on A cannot.)

The honest converse is our own boundary. At the single easy hop where a wrapper \emph{does} apply,
handing the verdict back in-loop is mildly \emph{harmful} ($\Delta$B, Appendix~D) --- so we do
not rest the headline payoff on the single-agent case. The concurrency case without a transactional
store is where the entire \emph{incremental} payoff over a wrapper lives.

\subsection{Trained classifiers}\label{sec:positioning-classifier}

The nearest neighbour is a \emph{trained} classifier: a model fitted on a labelled set of past failures, which then predicts whether a new run failed (arXiv 2511.04032\cite{silentfailures}, 2025, about 98\% accuracy on its own data). On paper that beats our detectors by a wide margin. The catch is what the number costs, and where that 98\% actually lives.

A classifier needs a labelled set of failures from the same kind of task, and a held-out slice to score it honestly. That labelled set is exactly the thing you do not have on a task you have never seen, which is the whole problem. Run the classifier \emph{fairly} on this corpus and two things happen. Its best-balanced operating point just predicts ``everything failed'': on a benchmark where 76.2\% of runs fail, that scores well by accuracy but buys no real skill, and this is where the 98\% headline lives (the recall mirage from \cref{sec:baselines}). Its only deployable operating point --- the one that almost never cries wolf --- fires exactly once. So at a false-alarm rate you could ship, it does nothing.

\begin{table}[t]
  \centering
  \caption{\cref{tab:classifier}. The zero-training detectors versus a trained classifier on this corpus, on the same scoreboard. A classifier has to pick an \emph{operating point} (how eager to fire); we show its two honest choices. Watch the \emph{lift} column --- the real-skill number from \cref{sec:bench} --- not the precision.}
  \label{tab:classifier}
  \begin{tabular}{lrrrr}
    \toprule
    method & fire & precision & lift & false-alarm \\
    \midrule
    \code{terminal\_error} (0-training) & 1.2\% & 95.0\% & +18.8pp & 0.24\% \\
    DOS trio (union, 0-training) & 5.1\% & 92.6\% & +16.4pp & 1.6\% \\
    trained clf --- \textbf{deployable} point & 0.0\% & 100.0\% & n/a (1 fire) & 0.0\% \\
    trained clf --- \textbf{F1-optimal} point & 100.0\% & 76.2\% & +0.0pp & 99.9\% \\
    \bottomrule
  \end{tabular}
\end{table}

The two classifier rows are the ends of one dial. At the \emph{deployable} point (the only setting that does not constantly cry wolf) it fires once in 6,862 runs --- effectively silent. At the \emph{F1-optimal} point (best by the usual accuracy measures) it just predicts ``everything failed'': 100\% fire, +0.0 lift, no skill. The widely-quoted ``98\% accuracy'' lives at that second setting --- the recall mirage of \cref{sec:baselines} in classifier's clothes. No dial position both fires usefully and stays trustworthy on an unseen task; our zero-training detectors sit at +16 to +19 lift at a deployable false-alarm rate.

Two further problems. The classifier reads the agent's own trajectory structure, so it risks the mirror-verifier trap (\cref{sec:intro}) --- it can learn the tells a model stops emitting once trained against. And when we feed it the DOS detector flags as features, its top weights \emph{are} those detectors (\code{dangling\_fired}, \code{terminal\_error\_fired}, \code{tool\_stream\_fired}): it re-derives the cheap signals it was handed. The right place for such a model is not beside DOS as a rival but \emph{under} it --- a JUDGE-rung adjudicator (ORACLE $\rightarrow$ JUDGE $\rightarrow$ HUMAN) that rules only on the residue the deterministic detectors abstained on, fail-to-abstain, never the floor.

\subsection{Arbiter neighbours}\label{sec:positioning-arbiter}

Two multi-agent arbitration systems sit nearby: Limen\cite{limen} (advisory write-leases) and CodeCRDT\cite{codecrdt} (lock-free CRDT). Both compare to DOS's \code{arbitrate()} call, not its detector, and neither can run on this benchmark: a frozen single-agent trace has no live write to mediate. (The arbiter axis is not unmeasured in this paper --- \cref{sec:payoff-race} exercises it live and counts J~=~6; neither neighbour reports a comparable prevented-corruption count.) On this frozen target DOS's arbiter likewise has nothing to mediate, and its enforcement point --- a pre-tool \code{permissionDecision: deny} hook an agent runtime such as Claude~Code honours before a tool runs --- is never exercised by a single-agent benchmark. The detector axis of this study and the arbiter axis of these neighbours do not share a scoreboard.

\section{Limitations}\label{sec:limits}

\begin{itemize}
  \item \textbf{In-loop: DETECT, not FIX --- no task-lift number on Toolathlon.} Every Toolathlon figure is detection on frozen runs; the one validated \emph{in-loop} action is give-up (\cref{sec:giveup}), which repairs nothing. The \emph{out-of-loop} result (\cref{sec:payoff}, J~=~10 across two models / 120 tasks) is a measured \emph{existence} result at an 8.3\% rate, not a calibrated population rate, and on a different benchmark (tau2-bench).
  \item \textbf{The recall ceiling is hard and tightens on the frontier.} Union recall is 6.18\%--7.42\%, so $\sim$93\% of failures stay silent, and fire-rate falls toward zero on the strongest models.
  \item \textbf{Four syscalls are null by construction here.} Toolathlon runs one agent on isolated state with no git or WAL, so \code{arbitrate}, \code{verify}, \code{liveness}, and \code{resume} have no job; only the in-flight detectors do.
  \item \textbf{The give-up token-saving magnitude is single-corpus, single-model.} The gate never wrongly halts a winner where it fires (specificity, across 22 models and two benchmarks); how \emph{often} it fires, and so how much it saves, is benchmark-dependent.
\end{itemize}

\section{Conclusion}\label{sec:conclusion}

Hold one thing fixed --- a byte-clean verdict the agent cannot forge, scored by a third-party
oracle DOS does not own --- and the result is about \emph{where} to spend it. \emph{In-loop},
every active fix (warn, rewind, inject the environment's own correction) was flat or
net-negative, because injecting a turn perturbs runs that would have passed; the one in-loop action
that helped was the negative one, give up correctly (\textbf{zero of 1,634 winners halted across 22
models}, $\sim$11\% of tokens saved). \emph{Out-of-loop}, the same verdict pays on both ways a
fleet fails --- a write-admission gate blocked \textbf{10 over-claims} at an identical 8.3\%
rate across two models (so the signal does not fade at the stronger tier), and an arbiter prevented
\textbf{all 6 would-be clobbers across 8 racing pairs} --- both off a byte no agent wrote,
neither touching the judged agent. The decisive move is not a better detector but a better \emph{consumer}: detect cheaply
in-loop, stop the doomed run, and spend the verdict out-of-loop --- with the W3 residue (a goal
with no world footprint) still beyond any witness short of a human.

\subsection{Directions}\label{sec:conclusion-directions}

The result generalizes along the axis that gave it: \emph{where} a non-forgeable verdict is
spent, and \emph{what witness} it can read. Four directions follow from the honest edges above, each
a measurement we know how to run. \textbf{Close the RLVR loop to a weight update.} \cref{sec:payoff-scope} measures the
label purge (poison reward labels removed off ground truth); the next rung is the trained-behaviour
delta --- fine-tune two models on the poison and the witness-cleaned reward set and measure
whether the second over-claims less. The pipeline is built and the proxy separates (+60pp over a
control that over-claims 0/4); only the weight update is unmeasured, and that is the one number a
training lab would want first. \textbf{Raise the witness floor.} About 38\% of frontier
goals reach \emph{no} sound witness at any rung (\cref{sec:boundary}); the open question is
how far a richer, still-non-forgeable witness (a structured state-diff, an environment-authored
trace, a second independent grader) can push that floor down before the irreducible human-seed
residue. \textbf{Cascade depth as a calibrated curve, not an existence shape.}
Appendix~D shows the payoff compounds F$^{D}$ across a branching fleet, but
with an injected root and per-node re-inheritance; a fully serial chain whose nodes \emph{write}
on the poison turns the shape into a calibrated constant. \textbf{Push the referee from advisory to
enforcing, carefully.} Everything here is advisory --- it proposes, a human or a
downstream system decides --- which is the right default; the controlled question is where an
out-of-loop \emph{enforcement} point (a pre-write deny the runtime honours) can be trusted to act
without a human in the loop, and where it cannot. The throughline is a small trusted referee that lets
a fleet of fallible, self-narrating agents --- and the loop that trains them --- be scaled
without scaling the trust placed in any one of them.

\section{Reproducibility}

Every number in the Toolathlon study (the detectors and their ceiling, \cref{sec:bench}--\cref{sec:result}; the give-up specificity, \cref{sec:giveup}; and the naive baselines, \cref{sec:baselines}) recomputes offline, with no network or model calls, from the durable per-run join \code{replay\_all\_rows.csv} (7,116 records, content fingerprint \code{sha256:1a55e8d8e2d4} --- computed line-ending-normalized, CRLF$\rightarrow$LF, so a checkout on Windows whose raw \code{sha256sum} differs is \emph{not} seeing drift; \code{additivity.py:rows\_fingerprint} is the canonical computation). The additivity single-source-of-truth re-asserts the headline claims and exits non-zero if a regenerated file drifts from them, so a claim cannot silently part from its data; the give-up cross-benchmark numbers and figure regenerate from their own script. The one exception is \cref{sec:payoff}: it reports a \emph{live, paid} run (tau2-bench through Gemini; single-model run \$0.89, cross-model run \$4.77), so its numbers are a fixed read-off of that run, not an offline recompute. They are transcribed verbatim into the figure's single source of truth (\code{paper/\_VERIFIED\_FACTS\_228\_2026-06-08.md} and \code{\_VERIFIED\_FACTS\_232\_2026-06-08.md}), and the run is resumable from a per-task cache so it never re-spends; re-running it requires a Gemini API key and reproduces an \emph{existence} result (the gate blocks real over-claims), not the exact count, since the live rollouts vary.

\begin{verbatim}
# 1. regenerate the durable rows from the frozen trajectories (offline once _data/ is cached):
python -m benchmark.toolathlon.run_replay --all --no-download --by-model \
    --rows-out benchmark/toolathlon/_results/replay_all_rows

# 2. recompute + re-assert the detector / additivity / union claims (exits non-zero on drift):
python -m benchmark.toolathlon.additivity --check

# 3. regenerate the give-up-correctly numbers and the cross-benchmark figure:
python benchmark/giveup_cross_benchmark.py

# 4. redraw every figure from the same single-source-of-truth numbers:
python -m benchmark.toolathlon.viz
\end{verbatim}

\textbf{Artifact.} The replay harness, the per-run data, the figures, and the single-source-of-truth scripts that compute every number in this paper are released together, under \code{benchmark/toolathlon/} and \code{benchmark/} in the project repository. The Toolathlon benchmark and its \code{Toolathlon-Trajectories} dataset (CC-BY-4.0) are the work of the HKUST-NLP authors; we replay them here as an external, third-party-scored test target, and they are not a contribution of this paper.
